% ===================================================================
%  அட்டவணை எண் 9 — மலை. நீரூற்று நகரம். காடு. வேட்டை.
%
%  Traduction, faite en 2026, en tamoul d'aujourd'hui, sur l'ido de
%  texto/io. Regles de la colonne : voir 10-attavanai-01.tex. Deux
%  scenes, deux numerotations : les cles portent c1 et c2.
%
%  TROISIEME COLONNE SUR CETTE PLANCHE, ET LE COMPTE SUIT LA
%  LATITUDE. Le coreen y perdait quatre noms — le chamois, la
%  bruyere, la marmotte, la pervenche. Le telougou en empruntait
%  neuf : hetre, chene, orme, peuplier, sapin, bruyere, fougere,
%  primevere, pie. Le tamoul en emprunte dix-sept : ஃபிர் le sapin,
%  சாமுவா le chamois, மார்மொட் la marmotte, ஹீதர் la bruyere,
%  ப்ரிம்ரோஸ் la primevere, பெரிவிங்கிள் la pervenche, ஸ்ட்ராபெரி la
%  fraise des bois, ட்ரஃபிள் la truffe, பீச் le hetre, பிர்ச் le
%  bouleau, ஓக் le chene, ஐவி le lierre, எல்ம் l'orme, பாப்ளர் le
%  peuplier, மேக்பை la pie, ரோ மான் le chevreuil, பார்ஃபிரி le
%  porphyre.
%
%  Coree 33-43° nord, quatre absents. Andhra 13-19° nord, neuf.
%  Pays tamoul 8-13° nord, dix-sept. Trois colonnes, trois
%  latitudes, et la courbe est monotone. Ce n'est plus une
%  illustration de la regle « le mot est la ou la chose est » : c'est
%  une mesure, et elle ne pouvait pas se faire avec moins de trois
%  colonnes sur la meme planche.
%
%  MAIS LE TAMOUL GARDE CE QUI VIT CHEZ LUI, ET IL EN GARDE
%  BEAUCOUP : எரிமலை le volcan — mot tamoul de fond, « la montagne de
%  feu » —, கணவாய் le defile, அருவி la cascade, ஏரி le lac, குகை la
%  grotte, கழுகு l'aigle, கரடி l'ours, சிலந்தி l'araignee, நத்தை
%  l'escargot, எறும்புப் புற்று la fourmiliere, விரியன் la vipere,
%  சாரைப் பாம்பு la couleuvre, காளான் le champignon, மரங்கொத்தி le
%  pic-vert, அணில் l'ecureuil, அன்னம் le cygne, பாசி la mousse,
%  பிசின் la resine, தேவதாரு le cedre. Et பெரணி la fougere, que le
%  telougou avait du emprunter : elle pousse dans les Ghats.
%
%  ET LA LIEUE NE TOMBE PAS JUSTE, POUR LA PREMIERE FOIS. Le releve
%  avait note deux coincidences : l'ఆమడ telougou et le 십 리 coreen
%  font tous deux quatre kilometres, exactement le « quaropa
%  kilometri » du bloc c1-15-1. J'en avais tire, dans la section
%  coreenne de LISEZ-MOI, qu'une journee de marche d'homme ne depend
%  d'aucune langue. C'est faux, et le tamoul le montre : son unite de
%  longue distance est le காதம், qui vaut selon les textes dix a
%  seize kilometres — pas quatre. La colonne ecrit donc நான்கு
%  கிலோமீட்டர். Deux coincidences ne font pas une loi ; la troisieme
%  langue etait celle qui pouvait la casser, et elle l'a cassee.
%
%  கோட்டை EST ECRIT ICI, ET C'EST LE CONTRAIRE DU TABLEAU 4. Au
%  tableau 4 j'avais refuse கோட்டை pour le « kastelo » (1), un manoir
%  de campagne, parce que le mot fait une forteresse. Le renvoi (8)
%  de ce tableau EST une forteresse — celle qui defend le defile
%  entre France et Espagne. Le mot juste au tableau 9 est le mot faux
%  au tableau 4 : ce n'est pas le mot qu'on juge, c'est la chose.
%
%  LE « cigano » (41), CINQUIEME FOIS DU RELEVE et deuxieme fois de
%  cette colonne apres le tableau 3. Le livret designe un homme par
%  son peuple ; la colonne ecrit ஊர் சுற்றும் மனிதர். La scene ne perd
%  rien : ce qui la fait tenir est le contraste du bonnet de fourrure
%  et du capuchon dechire, et il est intact.
%
%  NEUF ECARTS AU PREMIER PASSAGE, LE PIRE COMPTE DE LA COLONNE, ET
%  IL FAUT DIRE POURQUOI. Ce tableau est le premier ou j'ai ecrit
%  vite en me fiant a la liste de parigi.py — et sept des neuf ecarts
%  portent sur des couples que la liste NE CONTIENT PAS : la scierie
%  et le bucheron, le cor et les veneurs, la corneille et la
%  branche, l'araignee et sa mouche. Ce sont des ordres internes de
%  phrase, pas des rapports modificateur–tete. La lecon du tableau 5
%  se repete a une echelle plus grande : la liste dit ce qu'il faut
%  PREVOIR, elle ne dit pas ce qu'on a ECRIT, et se fier a elle seule
%  fait ecrire plus vite et plus faux.
%
%  ET J'AVAIS INVENTE UNE CLE, « t09-c1-05-1b ». Le fac-simile donne
%  deux alineas sous la seule cle c1-05-1, le second sans %%K parce
%  qu'il continue le meme numero. J'ai cru devoir en faire un bloc.
%  renvoji.py a repondu « cle inconnue » et kolonoj.py a signale le
%  bloc tronque a 29 % de la longueur mediane. Les cles sont celles
%  de texto/io, macro pour macro : on n'en ajoute pas, meme pour
%  ranger. Deux paragraphes dans un bloc s'ecrivent avec une ligne
%  blanche, comme au bloc t01-09-2.
%
%  TRENTE-CINQ PAIRES, DIX-HUIT RETOURNEMENTS. Le tableau decrit un
%  paysage et le paysage ATTACHE — la crete DE la chaine, le flanc DE
%  la montagne, l'aigle QUI a son nid, l'araignee QUI tisse sa toile.
%  Mais il enumere aussi beaucoup, arbres et betes, et l'enumeration
%  ne coute rien : c'est pourquoi le taux tient le milieu.
% ===================================================================

%%K t09-apar-1 apar
\VUsaut{4.0mm}
\VUtitre{15.0pt}{60}{அட்டவணை எண் 9}
\VUsaut{3.0mm}

%%K t09-tit-1 sub
\VUcentre{12.0pt}{0}{\textit{முதல் காட்சி.}}
\VUsaut{3.0mm}

%%K t09-tit-2 sub
\VUpk{11.4pt}{40}{மலை. --- நீரூற்று நகரம்.}
\VUsaut{3.0mm}

%%K t09-c1-01-1 p
1. --- எட்டு நாள்களாக நான் ப்ராட்ஸில் இருக்கிறேன். பிரெனேயின் வியத்தகு
\VUgras{மலைத் தொடருக்கு} \textsuperscript{(1)} முன்னால் நின்ற போது நான்
உணர்ந்த வியப்பு முழுவதையும் சொல்ல முடியவில்லை ; அதன் \VUgras{முகடு}
\textsuperscript{(2)} நீல வானத்தின் மேல் மாபெரும் தோரணம் போலத்
தெரிகிறது.

%%K t09-c1-02-1 p
2. --- பிரெனேயின் \VUgras{சிகரங்கள்} \textsuperscript{(3)} இவ்வளவு பெரிய
\VUgras{உயரத்தை} அடையும் என்று நான் நினைக்கவில்லை ; மிக அழகான
\VUgras{பனியாறுகளுக்கும்} \textsuperscript{(5)} பனி மூடிய \VUgras{மலை
உச்சிகளுக்கும்} \textsuperscript{(4)} முன்னால் நான் வியந்து
நின்றேன் ; அந்த உச்சிகளின் மேல் அச்சுறுத்தும் \VUgras{பனிச்
சரிவுகள்} உருள்கின்றன.

%%K t09-tit-3 sub
\VUsaut{3.0mm}
\VUpk{10.2pt}{40}{மலை நிலக்காட்சி.}
\VUsaut{3.0mm}

%%K t09-c1-03-1 p
3. --- நான் வந்த மறு நாளே, ப்ராட்ஸுக்கு அருகில் உள்ள அணைந்து போன பழைய
\VUgras{எரிமலைக்குப்} \textsuperscript{(6)} போனேன். \VUgras{எரிமலை வாயின்}
வறண்ட, மொட்டையான விளிம்புகளில், பல நூற்றாண்டுகளுக்கு முன்
\VUgras{மலைச் சரிவின்} \textsuperscript{(7)} மேல் ஓடிய \VUgras{எரிமலைக்
குழம்பின்} தடங்களை இன்னும் நன்றாகப் பார்க்க முடிகிறது.
\VUgras{சரிவுகளுள்} ஒன்றின் மேல் அந்தக் குழம்பின் சில துண்டுகளை
என்னால் கண்டுபிடிக்க முடிந்தது.

%%K t09-c1-04-1 p
4. --- ப்ராட்ஸ் எல்லையில் அமைந்திருக்கிறது. அங்குள்ள
\VUgras{கோட்டை} \textsuperscript{(8)} — அது பிரான்சையும் ஸ்பெயினையும்
பிரிக்கும் \VUgras{கணவாயைக்} \textsuperscript{(27)} காக்கிறது —,
ரைன் நதிக் கரையின் அந்தப் பழைய
கோட்டைகளை முற்றிலும் நினைவூட்டுகிறது ; அவை தங்கள் பாறையின் உச்சியிலிருந்து
\VUgras{இரையைப்} பார்த்துக் கொண்டிருப்பது போலத் தெரியும்.

%%K t09-c1-05-1 p
5. --- மாபெரும் \VUgras{கழுகு} \textsuperscript{(9)} அடிக்கடி என்
கவனத்தை ஈர்க்கிறது ; அது \VUgras{கோட்டைக்கு} \textsuperscript{(8)}
வெகு தொலைவில் இல்லாமல் கூடு கட்டியிருக்கிறது.

அந்த \VUgras{இரை கொல்லும் பறவைக்கு} வலிமையான \VUgras{அலகு} உண்டு. தன்
\VUgras{நகங்களால்} பிடித்து ஆட்டுக் குட்டியையோ வெள்ளாட்டுக் குட்டியையோ
அது அடிக்கடி தன் குஞ்சுகளுக்குக் கொண்டு வருகிறது. அதன்
\VUgras{சிறகுகள்} அவ்வளவு பெரியவை ; அதனால் அந்தப் பாரமான சுமையுடனும்
நெடு நேரம் அது சறுக்கிப் பறக்க முடியும்.

%%K t09-tit-4 sub
\VUsaut{3.0mm}
\VUpk{10.2pt}{40}{உலாச் செல்பவன்.}
\VUsaut{3.0mm}

%%K t09-c1-06-1 p
6. --- அங்கே மிகவும் இனிமையான உலா « கௌ » என்னும் அருவிக்குச்
\textsuperscript{(10)} செல்வதுதான். \VUgras{நீர் வீழ்ச்சி} சுழித்து
விழுந்து, பிறகு அழகான \VUgras{ஓடையாக} மறைந்து விடுகிறது ; அது
நகரத்திற்கு மேற்கே அமைந்த \VUgras{ஃபிர் மரக் காட்டில்}
\textsuperscript{(11)} மறைகிறது. இந்தக் காட்டின் வழியாக \VUgras{வானிலை
ஆய்வகம்} \textsuperscript{(12)} வரை ஏறினோம் ; \VUgras{காட்சி மேடையின்}
உச்சியிலிருந்து அந்தப் பகுதியின் \VUgras{பரந்த காட்சி} முழுவதையும்
பார்த்தோம். \VUgras{நிலக்காட்சி} வியத்தகு அழகு ; \VUgras{இயற்கை}
தன்னிடம் உள்ள எல்லாச் செல்வங்களையும் இந்த நாட்டுக்கு வாரி
வழங்குவது போலத் தெரிகிறது.

%%K t09-c1-07-1 p
7. --- இறங்குவதற்கு \VUgras{கயிற்று ரயிலை} \textsuperscript{(13)}
உபயோகித்தோம் ; ஒரு சிறு \VUgras{நிலையத்தில்} நின்றோம் ; அது
ப்ராட்ஸின் பிரபுக்கள் ஒரு காலத்தில் வாழ்ந்த பழைய \VUgras{நிலப்பிரபு
மாளிகையின்} \VUgras{இடிபாடுகளுக்கு} \textsuperscript{(14)} அருகில்
இருக்கிறது.

%%K t09-c1-08-1 p
8. --- பரிதாபமான மாளிகை! இப்போது நாகரிகமான \VUgras{வெந்நீரூற்று
நிலையமாக} இருக்கும் அந்த அழகான \VUgras{நீரூற்று நகரத்தைக்}
\textsuperscript{(15)} கீழே பார்க்கும் போது அது என்ன நினைக்கிறதோ?

%%K t09-c1-08-2 p
\VUgras{காலநிலை} அவ்வளவு இனிமையானது, \VUgras{கனிம நீர்} அவ்வளவு
பயன் தருவது ; அதனால் \VUgras{நலவாழ்வு இல்லத்தில்} \textsuperscript{(17)}
செய்யும் \VUgras{சிகிச்சை} உண்மையான மகிழ்ச்சியாக இருக்கிறது.

%%K t09-tit-5 sub
\VUsaut{3.0mm}
\VUpk{10.2pt}{40}{இனிமையான வெந்நீரூற்று நகரம்.}
\VUsaut{3.0mm}

%%K t09-c1-09-1 p
9. --- மேலும், தங்கியிருப்பதை இனிமையாக்க எல்லாம் செய்திருக்கிறார்கள்.
இரயில் பாதைக்கு வழி செய்ய பெரிய \VUgras{பாலப் பாதை} \textsuperscript{(16)}
கட்டப்பட்டது ; \VUgras{வெந்நீரூற்று நிலையத்தின்} \VUgras{பூங்கா}
\textsuperscript{(18)} — அங்கே \VUgras{இசை நிகழ்ச்சிகள்} நடக்கின்றன —
அழகான முறையில் அமைக்கப்பட்டிருக்கிறது. நேர்த்தியான « \VUgras{தனி மாளிகைகள்} »
\textsuperscript{(19)} பல \VUgras{சூதாட்ட விடுதியைச்}
\textsuperscript{(21)} சுற்றிக் கட்டப்பட்டன ; நன்றாகப் பராமரிக்கப்பட்ட
\VUgras{நெடுஞ்சாலை} \textsuperscript{(22)} போக்குவரத்தை மிகவும்
எளிதாக்குகிறது.

%%K t09-c1-10-1 p
10. --- எளிய \VUgras{சரிவு} \textsuperscript{(23)} \VUgras{சாலையை}
\textsuperscript{(20)} உண்மையான \VUgras{பள்ளத்தாக்கிலிருந்து}
\textsuperscript{(25)} பிரிக்கிறது ; அந்தப் பள்ளத்தாக்கின் அடியில்
அழகான சிறிய \VUgras{ஏரி} \textsuperscript{(26)} கிடக்கிறது ; அது
தெளிந்த \VUgras{ஓடைக்கு} \textsuperscript{(24)} நீர் தருகிறது.

%%K t09-c1-11-1 p
11. --- பள்ளத்தாக்கின் மறு பக்கத்தில் \VUgras{இரும்புச் சுரங்கம்}
\textsuperscript{(28)} தோண்டப்படுகிறது. \VUgras{சுரங்கக் குழியில்}
இறங்கும் \VUgras{சுரங்கத் தொழிலாளர்களைப்} பார்க்க நான் பல முறை
போனேன் ; அவர்கள் தங்கள் நாளைக் கழிக்கும் \VUgras{சுரங்கப் பாதைக்குள்}
கூட நுழைந்தேன் ; அங்கே \VUgras{உலோகம்} கொண்ட \VUgras{தாதுவை}
வெட்டி எடுக்கிறார்கள்.

%%K t09-c1-12-1 p
12. --- சுரங்கத்திலிருந்து எடுக்கும் செல்வங்களை உடனே பயன்படுத்த,
சுரங்கத்திற்கு அருகில் பெரிய \VUgras{உருக்காலை}
கட்டப்பட்டிருக்கிறது. இங்கே ஏராளமாகக் கிடைக்கும் \VUgras{நிலக்கரியால்}
சூடாக்கப்படும் \VUgras{பெரும் உலை} \textsuperscript{(29)}
\VUgras{வார்ப்பிரும்பை} விரைவாகவும் மலிவாகவும் பெற வழி செய்கிறது.

%%K t09-c1-13-1 p
13. --- சுரங்கத்திற்கு வெகு தொலைவில் இல்லாமல் \VUgras{கல் குவாரி}
\textsuperscript{(30)} தோண்டப்படுகிறது. பழைய \VUgras{கல் வெட்டுபவர்}
தன் \VUgras{கடப்பாரையால்} \VUgras{கருங்கல்லையோ},
\VUgras{பார்ஃபிரியையோ}, \VUgras{மணற்கல்லையோகூட} வெட்டுவதில்லை ;
வீடு கட்டுவதற்கான எளிய \VUgras{செதுக்கிய கல்லையே} வெட்டுகிறார்.

%%K t09-c1-14-1 p
14. --- அந்த நாட்டின் மிக அழகான அலங்காரங்களுள் ஒன்று \VUgras{ஃபிர்
மரங்கள்} \textsuperscript{(31)}. \VUgras{பாறை இடுக்கின்}
\textsuperscript{(32)} அடியில், \VUgras{வேகமான ஓடையின்}
\textsuperscript{(33)} கரையில், \VUgras{படுகுழியின்}
\textsuperscript{(34)} அல்லது செங்குத்துப் பாறையின் விளிம்பில் —
எங்கும் அவற்றின் மெல்லிய \VUgras{ஊசி இலைகள்} பசுமையாக மின்னுகின்றன.

%%K t09-tit-6 sub
\VUsaut{3.0mm}
\VUpk{10.2pt}{40}{நடைப் பயணம்.}
\VUsaut{3.0mm}

%%K t09-c1-15-1 p
15. --- \VUgras{நெடு நேரம்} உலாவுவதற்கு என் விடுமுறையைப்
பயன்படுத்துகிறேன். மலையின் மேல் வளைந்து செல்லும் \VUgras{பாதைகள்}
\textsuperscript{(35)}, தூரத்தைப் பொருட்படுத்தாமல் பல
\VUgras{கிலோமீட்டர்} — அடிக்கடி பல \VUgras{நான்கு கிலோமீட்டர்}
தொலைவுகள் — கடக்க வழி செய்கின்றன.

%%K t09-c1-15-2 p
\VUgras{எல்லைக் கற்களும்} \textsuperscript{(36)} \VUgras{வழிகாட்டித்
தூண்களும்} \textsuperscript{(37)} பாதைகளைக் குறிக்கின்றன ; அவற்றின் மேல்
அடிக்கடி இளம் \VUgras{மலைவாசியைச்} \textsuperscript{(38)}
சந்திக்கிறோம் — அவனுடைய பக்கத்தில் \VUgras{இரட்டைப் பை}
\textsuperscript{(40)} தொங்குகிறது, அவன் \VUgras{மார்மொட்டைச்}
\textsuperscript{(39)} சுமந்து வருகிறான் ; அல்லது \VUgras{ஊர் சுற்றும்
மனிதரைச்} \textsuperscript{(41)} சந்திக்கிறோம் — அவருடைய \VUgras{ரோமத்
தொப்பி} \textsuperscript{(42)} கிழிந்த பழைய \VUgras{முக்காட்டுடன்}
\textsuperscript{(43)} வித்தியாசமாக முரண்படுகிறது. அவர் பழக்கிய
\VUgras{கரடியைச்} \textsuperscript{(44)} சங்கிலியால் இழுத்து
வருகிறார்.

%%K t09-tit-7 sub
\VUpk{10.2pt}{40}{மலை ஏறுதல்.}
\VUsaut{3.0mm}

%%K t09-c1-16-1 p
16. --- கிட்டத்தட்ட நாள்தோறும் \VUgras{வழிகாட்டியுடன்}
\textsuperscript{(48)} \VUgras{ஏற்றம்} பயில்கிறேன் ; முதுகில்
\VUgras{பையுடன்} \textsuperscript{(47)}, கையில் « \VUgras{மலை
ஊன்றுகோலுடன்} », உண்மையான \VUgras{சுற்றுலாப் பயணியைப்}
\textsuperscript{(46)} போல \VUgras{பாறைகளைத்} \textsuperscript{(45)}
தாக்கும் போது நான் மிகவும் பெருமிதம் கொள்கிறேன்.

%%K t09-c1-17-1 p
17. --- மலையின் உச்சியிலிருந்து பார்த்தால், சமவெளியின் மக்கள்
எறும்புகளை விடப் பெரியவர்களாகத் தெரியவில்லை ; \VUgras{செம்மறி மந்தை}
\textsuperscript{(50)} — அது \VUgras{மேய்ச்சல் நிலத்தில்}
\textsuperscript{(49)} மேய்கிறது — குழந்தைகளுக்கான விளையாட்டுப்
பொருள் போலத் தெரிகிறது. \VUgras{பைன் மரமும்} \textsuperscript{(51)}
\VUgras{மரச் சிறு வீடும்} \textsuperscript{(52)} பசுமையின் மேல் ஒரு
புள்ளியாகவே தெரிகின்றன.

%%K t09-c1-18-1 p
18. --- இந்த உலாக்களின் போது \VUgras{ஒற்றையடிப் பாதையில்}
\textsuperscript{(53)} \VUgras{சாமுவா வேட்டைக்காரனை}
\textsuperscript{(54)} அடிக்கடி சந்திக்கிறோம் ; அவன் இப்போதுதான்
கொன்ற விலங்கைத் தோளில் சுமந்து வருகிறான்.

%%K t09-c1-19-1 p
19. --- என் தாவரத் தொகுப்பில் மிகக் குறிப்பிடத்தக்க \VUgras{பெரணிக்}
\textsuperscript{(55)} கிளை ஒன்றையும், அழகான ரோஜா நிற \VUgras{ஹீதர்}
\textsuperscript{(57)} கிளை ஒன்றையும் வைத்தேன் ; அவற்றை என் கடைசி
உலாவின் போது சிறிய \VUgras{குகையின்} \textsuperscript{(56)} வாசலில்
பறித்தேன்.

%%K t09-tit-8 sub
\VUsaut{3.0mm}
\VUcentre{12.0pt}{0}{\textit{இரண்டாம் காட்சி.}}
\VUsaut{3.0mm}

%%K t09-tit-9 sub
\VUpk{11.4pt}{40}{காடு. --- வேட்டை.}
\VUsaut{3.0mm}

%%K t09-c2-01-1 p
1. --- நேற்று காட்டில் இனிமையான ஒரு நாளைக் கழித்தேன். அங்கே வாழும்
விலங்குகளுக்கும் \VUgras{பூச்சிகளுக்கும்} \textsuperscript{(5)} இடையே
நிலவும் சுறுசுறுப்பு என்னை மிகவும் கவர்ந்தது.

%%K t09-c2-01-2 p
\VUgras{சிலந்தி} \textsuperscript{(1)} — அது இரண்டு கிளைகளுக்கு இடையே
தன் \VUgras{வலையை} \textsuperscript{(2)} நெய்கிறது, கடந்து போகும்
\VUgras{ஈயைப்} \textsuperscript{(3)} பிடிப்பதற்காக — ; தன் ஓட்டை மிகுந்த சிரமத்துடன்
சுமக்கும் \VUgras{நத்தை} \textsuperscript{(4)} ; தன் இரையைத் தேடும்
கொம்பன் வண்டு ; \VUgras{எறும்புப் புற்றுக்கு} \textsuperscript{(6)}
அருகில் மிகவும் ஆர்வமான உழைப்பாளி எறும்பு — இந்தச் சிறியவை
எல்லாம் போயின, வந்தன, அசாதாரண விடாமுயற்சியுடன் உழைத்தன.

%%K t09-c2-02-1 p
2. --- ஒரு \VUgras{பாம்பு} \textsuperscript{(7)} — முதல் பார்வையில் அதை
\VUgras{விரியன்} என்று நினைத்தேன், ஆனால் அது எளிய \VUgras{சாரைப்
பாம்பே} — மர அடிமரத்திற்குக் கீழே ஒளிந்திருந்தது ; அவ்வப்போது தன்
மெல்லிய சிறு தலையை வெளியே நீட்டியது ; அது எப்போதும் கடிக்கத் தயாராக
இருப்பது போலத் தெரிந்தது. என் முன்னால் ஒரு பருமனான \VUgras{நிலநத்தை}
\textsuperscript{(8)} — அங்கே ஏராளமாக வளரும் சுவையான \VUgras{ஸ்ட்ராபெரிப்
பழங்களுள்} \textsuperscript{(11)} ஒன்றை விழுங்கிய பிறகு — பாரமாக
ஊர்ந்தது.

%%K t09-c2-03-1 p
3. --- தரை \VUgras{காட்டு மலர்களால்} மூடப்பட்டிருந்தது :
\VUgras{ப்ரிம்ரோஸ் மலர்களும்} \textsuperscript{(9)},
\VUgras{பெரிவிங்கிள் மலர்களும்} \textsuperscript{(10)}, எனக்குத்
தெரியாத வேறு பல செடிகளும் \VUgras{புல் புதர்களுக்கு}
\textsuperscript{(12)} இடையே பூத்திருந்தன.

%%K t09-c2-04-1 p
4. --- அந்தப் பகுதியில் \VUgras{ட்ரஃபிள்} \VUgras{காளானைப்}
\textsuperscript{(14)} போலவே கிட்டத்தட்ட வழக்கமானது ; காட்டுக்
காவலருக்குச் சொந்தமான ஒரு \VUgras{பன்றிக் குட்டி}
\textsuperscript{(13)} அவற்றுள் சிலவற்றைக் கண்டுபிடிக்க முடியுமா
என்று பேராசையுடன் தேடியது.

%%K t09-tit-10 sub
\VUsaut{3.0mm}
\VUpk{10.2pt}{40}{திருட்டு வேட்டை.}
\VUsaut{3.0mm}

%%K t09-c2-05-1 p
5. --- என்னை மிகவும் மகிழ்வித்த ஒரு காட்சியின் \VUgras{முடிவை}
நான் பார்த்தேன். தன் \VUgras{குட்டைக் கால் நாயின்}
\textsuperscript{(15)} மோப்பத்தின் துணையால் \VUgras{காட்டுக் காவலர்}
\textsuperscript{(16)} இப்போதுதான் ஒரு \VUgras{திருட்டு
வேட்டைக்காரனைப்} \textsuperscript{(22)} பிடித்து விட்டார் ; அவன் தான்
கொன்ற \VUgras{வேட்டைப் பொருளைச்} \textsuperscript{(17)} சுமந்து செல்லத்
தயாராகும் தருணத்தில். கைது ஆவதை விடத் தன் வேட்டைப் பொருளை விட்டு
விடுவதையே விரும்பி அந்தத் திருட்டு வேட்டைக்காரன் ஓடி விட்டான் ;
\VUgras{முயலும்} \textsuperscript{(18)}, \VUgras{பெண் மானும்}
\textsuperscript{(19)}, \VUgras{ரோ மானும்} \textsuperscript{(20)},
\VUgras{காட்டுப் பன்றியும்} \textsuperscript{(21)} புல்லின் மேல்
கிடந்து காட்டுக் காவலரை மகிழ்வித்தன.

%%K t09-c2-06-1 p
6. --- காட்டில் உள்ள மரங்களின் வேறுபாடு வியப்பாக இருக்கிறது.
\VUgras{பீச் மரங்களும்} \textsuperscript{(23)}, \VUgras{பிர்ச்
மரங்களும்} \textsuperscript{(26)}, \VUgras{பிசினைத்} \textsuperscript{(27)}
தரும் \VUgras{பைன் மரங்களும்}, \VUgras{ஓக் மரங்களும்}
\textsuperscript{(29)}, வேறு பலவும் தங்கள் கிளைகளைக் கலக்கின்றன.

%%K t09-c2-06-2 p
உயரமான மரங்களின் அடிமரத்தில் ஒட்டிக் கொள்ளும் \VUgras{ஐவி கொடி}
\textsuperscript{(24)}, ஒளிரும் பசுமையால் \VUgras{காட்டுக்கு}
\textsuperscript{(25)} நிறம் தருகிறது ; அது \VUgras{பாசியின்}
\textsuperscript{(31)} இன்னும் தெளிவான, வெளிரிய பச்சை நிறத்துடன்
இனிமையாக இணைகிறது ; அந்தப் பாசியில் \VUgras{மரங்கொத்தி}
\textsuperscript{(30)} பூச்சிகளைத் தேடுகிறது.

%%K t09-tit-11 sub
\VUsaut{3.0mm}
\VUpk{10.2pt}{40}{காட்டு நிலக்காட்சி.}
\VUsaut{3.0mm}

%%K t09-c2-07-1 p
7. --- பழைய ஓக் மரங்கள் என்னைக் கவர்ந்தன. நிலத்திலிருந்து பாதி
வெளிவந்த அவற்றின் பருமனான, முறுக்கிய \VUgras{வேர்கள்}
\textsuperscript{(32)} அவற்றுக்கு வலிமையின், வீரியத்தின் அற்புதமான
தோற்றத்தைத் தருகின்றன. \VUgras{சொந்தக்காரர்} \textsuperscript{(35)}
எப்படி அவற்றை வெட்டி \VUgras{அறுவாளுக்கு} \textsuperscript{(34)} —
\VUgras{மரம் வெட்டுபவரின்} \textsuperscript{(33)} அறுவாளுக்கு —
ஒப்படைக்க மனம் ஒப்புகிறார் என்று என்னையே கேட்டுக் கொள்கிறேன்.
பரிதாபமான \VUgras{அடிமரங்கள்} \textsuperscript{(36)} என்னைத்
துயரப்படுத்துகின்றன ; அவை தங்கள் \VUgras{பட்டையை}
\textsuperscript{(37)} இழந்தவை, அல்லது \VUgras{கோடரியால்}
\textsuperscript{(38)} பிளந்தவை — \VUgras{நாள்கூலி வேலையாளின்}
\textsuperscript{(39)} கோடரியால்.

%%K t09-c2-08-1 p
8. --- அருகில், பழைய \VUgras{கரி சுடுபவர்} \textsuperscript{(41)} தன்
\VUgras{மண் அள்ளியால் கரிக் குவியல்} \textsuperscript{(42)} ஒன்றைத்
தயார் செய்தார் ; ஒரு கிழவி வெட்டிய மரங்களின் சிறு கிளைகளால் ஆன
\VUgras{சுள்ளி விறகுக்} \VUgras{கட்டு} \textsuperscript{(40)} ஒன்றைச்
சுமந்து வந்தாள்.

%%K t09-tit-12 sub
\VUsaut{3.0mm}
\VUpk{10.2pt}{40}{வேட்டை.}
\VUsaut{3.0mm}

%%K t09-c2-09-1 p
9. --- சில கணங்கள் ஓய்வு எடுக்க \VUgras{காவலர் வீட்டுக்குப்}
\textsuperscript{(46)} போக விரும்பினேன் ; ஆனால் \VUgras{சிறு ஏரிக்கு}
\textsuperscript{(43)} அருகில் வந்த போது — அதன் மேல் அழகான
\VUgras{அன்னம்} \textsuperscript{(45)} நீந்துகிறது — சமீபத்திய
\VUgras{வெள்ளம்} \textsuperscript{(44)} அந்த வழியை உண்மையான
\VUgras{சதுப்பாக} மாற்றி விட்டதைக் கவனித்தேன். திரும்ப வேண்டியதாயிற்று.
காட்டின் விளிம்பில் இருக்கும் \VUgras{தேவதாருவுக்கும்}
\textsuperscript{(49)}, \VUgras{பாப்ளர் மரங்களுக்கும்}
\textsuperscript{(48)}, \VUgras{எல்ம் மரத்திற்கும்}
\textsuperscript{(47)} அருகில் கடந்து போகும் போது, ஒரு \VUgras{சிறு
காட்டிலிருந்து} \textsuperscript{(54)} மிக அழகான \VUgras{ஆண் மான்}
\textsuperscript{(50)} வெளிவருவதைப் பார்த்தேன் ; அதை விடாப்பிடியான
\VUgras{வேட்டை நாய்க் கூட்டம்} \textsuperscript{(51)} துரத்தியது ;
\VUgras{கொம்பு எக்காளமும்} \textsuperscript{(53)} அவற்றைத் தூண்டியது —
\VUgras{குதிரையில் வரும் வேட்டைக்காரர்களின்} \textsuperscript{(52)}
எக்காளம்.
பரிதாபமான அந்த விலங்கு முற்றிலும் சோர்ந்து போயிருந்தது. என்னிடமிருந்து
சில அடிகள் தொலைவில் அது சாய்ந்து இறந்தது.

%%K t09-c2-10-1 p
10. --- \VUgras{துரத்தும் வேட்டையின்} ஆரவாரம் ஈர்த்த காட்டுக்
காவலரின் மகன் வீட்டிலிருந்து வெளியே வந்தான் ; இருவரும் மீண்டும்
காட்டுக்குள் திரும்பினோம். அவன் பழைய ஓக் மரத்தின் கிளையில் ஒரு
\VUgras{மேக்பையைப்} \textsuperscript{(56)} பார்த்திருந்தான். அதன் கூடு
\VUgras{இலைத் தொகுதிக்குள்} \textsuperscript{(58)} மறைந்திருக்கும் என்று
அவன் நினைத்தான் ; அந்தக் கூட்டைப் பிடிக்க விரும்பினான்.

%%K t09-c2-10-2 p
\VUgras{அணிலைப்} \textsuperscript{(61)} போலச் சுறுசுறுப்பாக அவன்
\VUgras{கிளையிலிருந்து} \textsuperscript{(55)} கிளைக்கு ஏறினான் ; ஆனால்
எதுவும் கிடைக்கவில்லை ; பழைய \VUgras{காகத்தைப்} \textsuperscript{(60)}
பறக்க வைத்தது மட்டுமே அவனால் முடிந்தது — அது \VUgras{கிளை}
\textsuperscript{(57)} ஒன்றின் மேல் அமர்ந்திருந்தது.

\VUsaut{4.0mm}
\VUfilet{22mm}
